Volatility forecasting is a foundational input to financial risk management. Forward-looking volatility estimates are directly embedded in Value-at-Risk calculations, portfolio risk-limit setting, derivative pricing and dynamic hedging, margin and collateral protocols, and stress-test scenario calibration. The operational challenge is that daily stock-index volatility is generated by several concurrent mechanisms: strong short-run persistence, medium-horizon adjustment to macro-financial conditions, and episodic synchronized stress propagation across markets. A practical risk-monitoring system must therefore do more than fit average conditions. It must remain strictly causal to preserve the integrity of real-time risk assessments, beat credible persistence-based benchmarks to justify adoption costs, and continue to generate reliable warning signals precisely when market conditions become most unstable---and when the consequences of a missed or false alert are most severe \cite{corsi2008har,song2022garchmidas,zeng2024macrodriven}.

This requirement immediately creates a tension. On the one hand, parsimonious linear models such as HAR remain surprisingly strong in daily volatility forecasting, especially at the one-day horizon \cite{corsi2008har,branco2024linear}. On the other hand, a purely linear persistence view is too narrow to describe the full structure of market risk. Volatility clustering is not generated at a single frequency, implied-volatility measures and technical indicators interact nonlinearly, and stress episodes often unfold through cross-market transmission rather than through one index in isolation. The practical question is therefore not whether nonlinear models can be made more complex than HAR. It is whether a carefully designed multiscale representation can extract additional predictive information without sacrificing causality, interpretability, or reproducibility.

Three literature streams motivate this question. First, decomposition-based studies show that volatility prediction can benefit from multiscale analysis when signals evolve at heterogeneous frequencies \cite{souropanis2023wavelet,zhuo2024harsvr,ma2025autoformerapplsci,su2025hybridapplsci}. Second, machine-learning studies show that nonlinear learners can exploit richer public information sets, including implied volatility, technical signals, and macro-financial factors \cite{diaz2024mlvalue,zhang2023intradayml,chun2025timing}. Third, spillover studies document that market stress propagates across indices, although that propagation is usually modelled in the time domain rather than as explicit frequency-specific transmission \cite{son2023gnnspillover}. The literature on volatility warning and extreme-risk identification adds a fourth strand by showing that forecasting and alert generation should be connected more tightly than they often are in practice \cite{ren2024extreme,ran2024earlywarning,purnell2024explainable}.

Despite this progress, four gaps remain. First, many studies benchmark against weak baselines, which makes it difficult to determine whether the proposed method adds value beyond the persistence structure already captured by HAR. Second, wavelet features are usually extracted within one index only; relatively little evidence is available on whether peer-index wavelet components carry incremental predictive content at specific frequencies. Third, the regression task of volatility forecasting and the classification task of risk warning are commonly treated as separate pipelines even though they rely on the same chronology and largely the same public information. Fourth, reproducibility is often limited by proprietary data, ambiguous release timing for predictors, or full-sample preprocessing that quietly introduces look-ahead bias.

This paper addresses those gaps by proposing a public-data Causal Multiscale Wavelet Spillover Learning (CMWSL) framework for stock-index volatility forecasting and risk early warning. CMWSL is built around three linked components. The first is a HAR persistence block that preserves the strongest linear benchmark in the literature. The second is a causal stationary wavelet transform (SWT) representation that summarizes within-index short-, medium-, and longer-run dynamics using rolling scale-wise statistics. The third is a spillover layer that appends medium- and long-scale wavelet energy from peer indices in order to test whether cross-index information enters through specific frequency bands rather than only through contemporaneous correlations.

The central research question is deliberately narrow and testable: \emph{when} do multiscale features and peer-index wavelet spillovers add value beyond persistence in stock-index volatility forecasting? The paper does not seek to prove blanket dominance of a new model class. Instead, it asks whether the incremental value of multiscale learning is conditional on forecast horizon, target construction, information richness, and market state. This framing is more appropriate for risk management than a simple average-score horse race, because the operational value of a model often appears precisely when markets depart from their typical regime.

The article formalizes that question through the following hypothesis:

\begin{quote}
\textit{Cross-Index Wavelet Spillover Hypothesis.} Medium- and long-scale wavelet detail components (levels d2 and d3) of peer equity indices contain incremental predictive information for a target index's future average volatility, beyond the information contained in the target index's own persistence terms and within-index multiscale representation. This incremental value is state-dependent and horizon-dependent, and can be identified through nested out-of-sample forecast comparison.
\end{quote}

To test this hypothesis, the empirical design uses a three-step controlled ablation: \textit{HAR} as the persistence benchmark, \textit{wavelet\_lightgbm} as the within-index multiscale extension, and \textit{spillover\_lightgbm} as the same model further augmented with peer-index d2/d3 wavelet energy features. The study relies only on public daily data for the S\&P~500, Nasdaq-100, and Dow Jones Industrial Average, and evaluates all models under a 2016--2025 walk-forward protocol with strict causal feature construction. In addition to the regression task, the same pipeline is used to generate high-volatility warning labels and probability-based alerts, allowing forecast quality and warning quality to be studied within one coherent framework \cite{ren2024extreme,ran2024earlywarning,purnell2024explainable}.

\subsection{Related Literature and Research Positioning}

The paper sits at the intersection of four literature streams. First, persistence-based volatility models establish that any serious forecasting framework must compete against parsimonious linear benchmarks rather than circumvent them \cite{corsi2008har,song2022garchmidas,branco2024linear}. Second, wavelet and decomposition studies show that scale-specific summaries can improve volatility prediction when the data-generating process contains heterogeneous temporal frequencies \cite{souropanis2023wavelet,zhuo2024harsvr,ma2025autoformerapplsci,su2025hybridapplsci}, but most of this work remains within-market. Third, machine-learning and graph-based studies model cross-market dependence and nonlinear interaction structures \cite{diaz2024mlvalue,son2023gnnspillover,song2023cnn_gru,chun2025timing}, though mainly in the time domain rather than through explicit frequency-specific spillover channels. Fourth, early-warning studies translate continuous volatility and stress signals into actionable risk states \cite{ren2024extreme,ran2024earlywarning,purnell2024explainable}, but often as a downstream add-on rather than as part of the main forecasting design. Recent review evidence suggests that multi-frequency and cross-market synthesis remains a live research gap \cite{leushuis2026review}. Table~\ref{tab:lit_compare} positions the present study relative to these strands.

\begin{table}[htbp]
\centering
\footnotesize
\caption{Positioning of the present study relative to representative prior works.
  \checkmark\ $=$ present; $\circ$ $=$ partial; --- $=$ absent.
  \textit{Within-$\lambda$}: within-index wavelet decomposition;
  \textit{Cross-$\lambda$}: cross-index wavelet spillover features;
  \textit{Warning}: integrated early-warning layer;
  \textit{Public}: all data from public sources;
  \textit{Walk-fwd}: rolling walk-forward evaluation;
  \textit{DL base.}: deep-learning model included as baseline.}
\label{tab:lit_compare}
\begin{tabular}{lcccccc}
\toprule
Study & Within-$\lambda$ & Cross-$\lambda$ & Warning & Public & Walk-fwd & DL base. \\
\midrule
\citet{corsi2008har}          & ---         & ---         & ---         & ---         & $\circ$     & ---         \\
\citet{souropanis2023wavelet} & \checkmark  & ---         & ---         & $\circ$     & \checkmark  & ---         \\
\citet{zhuo2024harsvr}        & ---         & ---         & ---         & ---         & \checkmark  & ---         \\
\citet{diaz2024mlvalue}       & ---         & ---         & ---         & ---         & \checkmark  & \checkmark  \\
\citet{zhang2023intradayml}   & ---         & ---         & ---         & ---         & \checkmark  & \checkmark  \\
\citet{son2023gnnspillover}   & ---         & \checkmark  & ---         & ---         & ---         & \checkmark  \\
\citet{song2023cnn_gru}       & ---         & ---         & ---         & ---         & $\circ$     & \checkmark  \\
\citet{ran2024earlywarning}   & ---         & ---         & \checkmark  & ---         & ---         & ---         \\
\textbf{This study}           & \checkmark  & \checkmark  & \checkmark  & \checkmark  & \checkmark  & \checkmark  \\
\bottomrule
\end{tabular}
\end{table}

Relative to this literature, the contribution of the paper is fourfold.
\begin{enumerate}
\item It proposes an integrated causal pipeline for market risk management that links volatility forecasting and high-volatility early warning within a single public-data, walk-forward framework---eliminating the information gap that typically separates risk quantification from risk alert generation.
\item It formalizes cross-index wavelet spillover as a nested incremental-information problem, enabling a traceable assessment of which frequency-specific cross-market channels contribute to portfolio risk exposure beyond local persistence.
\item It evaluates model performance through risk-management-relevant diagnostics---tail-conditioned, market-state-conditioned, rolling-window, and stress-event-based decompositions---that reveal exactly when multiscale representation delivers actionable improvements over linear benchmarks.
\item It delivers the empirical study as reproducible risk-forecasting research built entirely on public market and macro-financial data, with deterministic scripts and reviewer-ready experiment artifacts consistent with a transparent, auditable risk-model evaluation standard.
\end{enumerate}

These contributions imply a deliberately disciplined empirical claim. The article does not argue that nonlinear multiscale learning should replace HAR everywhere. Instead, it argues that persistence remains the correct baseline at short horizons, while causal wavelet representation and peer-index spillover provide useful incremental information when the forecasting problem becomes less myopic, the information environment becomes richer, and market stress becomes more synchronized. That is precisely the regime in which practical risk monitoring is most demanding---and the setting to which this paper's contributions are directly addressed.

The remainder of the paper is organized as follows. Section~\ref{sec:data} describes the market sample, target construction, and predictor blocks. Section~\ref{sec:methodology} presents the CMWSL framework. Section~\ref{sec:design} details the walk-forward evaluation protocol and ablation design. Section~\ref{sec:results} reports the empirical evidence. Section~\ref{sec:discussion} interprets the findings, discusses risk management implications, and notes limitations; Section~\ref{sec:conclusions} concludes.
